\chapter{引言}\label{chap:introduction}

\section{研究背景与研究意义}

近年来，随着我国高等教育规模的不断扩张，多校区办学已成为高校优化办学空间与学科布局的常规模式。校区之间的通勤需求、校内外交通衔接需求以及节假日期间集中出行需求显著增加，使得传统校园内部交通系统面临着前所未有的承载压力。由于部分新校区地处城市边缘，外围公共交通的局限性进一步强化了师生对校内班车及自发拼车行为的依赖。在这种背景下，如何通过技术手段优化资源配置，已成为高校后勤治理中亟待解决的现实课题。

在现有的运行机制下，校园交通信息的高度碎片化构成了主要的治理瓶颈。班车排班信息通常散落于官网公告、社交群组或线下布告栏，缺乏统一的数字化入口；而拼车需求则高度依赖即时通讯工具等非正式渠道，呈现出极强的随机性与低透明度。这种现状导致师生在出行决策时面临较高的检索成本，也使拼车组织长期停留在低效率、弱约束的状态。

针对上述痛点，构建集约化、智能化的校内综合交通平台具有重要的现实意义。通过整合班车动态、拼车信息与用户交互服务，系统能够降低出行决策成本，提升供需对接效率，并借助校内身份认证与信用评价机制增强共享出行的安全性和可靠性。同时，该平台还能为班车线路优化与校园交通治理提供数据支持，具有明显的工程应用价值。

\section{国内外研究现状}

随着共享出行模式在城市交通体系中的广泛应用，围绕拼车调度与供需匹配的算法研究逐渐成为运筹优化与智能决策领域的重要方向。共享出行问题在理论层面可追溯至运筹学中的经典车辆路径问题（Vehicle Routing Problem, VRP）及其衍生出的拨号运输问题（Dial-a-Ride Problem, DARP）\cite{laporte1992,cordeau2007}。早期研究主要在静态需求下构建整数规划模型，通过分支定界\cite{huang2014}或基于拉格朗日列生成法（Lagrangean column generation）的精确算法寻求最优路径方案\cite{baldacci2004}。随着实时计算能力的提升，Agatz 等学者对动态拼车（Dynamic Ride-sharing）问题进行了系统综述，提出在时空约束条件下实现多对多匹配的优化框架\cite{agatz2012}。Yan 等学者采用时空网络流技术，构建多对多拼车问题的模型，涉及多种车辆类型和多人类型\cite{yan2011}。Luk 等学者则尝试利用图论工具分析拼车者特征以优化车辆分配\cite{knapen2014}。然而，在动态请求持续到达的环境中，精确算法往往面临计算复杂度迅速增长的问题，难以满足实时响应的需求。为解决这一问题，Najmi 等提出了动态实时拼车的匹配策略和目标函数，并采用滚动时域（rolling horizon）策略来动态匹配司机与乘客，同时关注实时性和规模性问题\cite{najmi2017}。此外，Schreieck 等专门研究了动态拼车中自动匹配算法的设计与实现，深入讨论了实时匹配的复杂性和算法策略\cite{schreieck2016}。

为了在解的质量与计算效率之间取得平衡，研究逐渐转向启发式与元启发式方法。Cordeau 等学者针对 DARP 模型提出改进的禁忌搜索算法，在大规模问题场景下显著提升了求解效率\cite{ho2018}。Huang 等学者开发了一款移动应用，基于真实用户数据的计算结果显示了该算法的有效性\cite{huang2016}。同时，遗传算法（GA）被广泛应用于动态出租车拼车模型\cite{huang2015}，而蚁群算法\cite{huang2019}、模拟退火算法\cite{jung2016}也能缓解高维组合优化的压力。这类方法虽然在一定程度上缓解了高维组合优化带来的计算压力，但其核心仍依赖预设的目标函数。在多目标条件下，路径偏移率、时间窗重合度及系统成本等因素通常需要通过人工设定权重进行综合评价，模型适配性与泛化能力存在一定局限。

近年来，随着数据获取能力的增强与机器学习技术的发展，研究者开始尝试引入数据驱动方法以刻画复杂的匹配关系。Wen 等提出了一种基于全球定位系统（GPS）的拼车路线匹配方案，用于城市拼车服务中的路线挖掘、乘客选择和路线合并\cite{he2014}。以集成学习为代表的算法在交通预测与出行决策领域表现更为突出，Breiman 提出的随机森林算法\cite{breiman2001}通过构建多棵决策树进行集成预测，具有较强的抗过拟合能力与稳定性。相关研究\cite{hagenauer2017,wu2024,anil2025}表明，该算法能够有效整合离散的时空约束与用户属性特征，并准确刻画个体在面对不同交通方案时的决策逻辑。因此，将拼车匹配过程中的时空约束与行为特征进行特征化表达，并通过随机森林模型进行匹配评分预测，在方法层面具有一定的理论依据。

随着我国高校多校区办学模式的普及，校区间通勤已成为高校管理和学生日常生活中的重要问题。现有研究表明，多校区办学模式下，校车运行成本显著增加，道路拥堵频发，且调度管理复杂\cite{wu2021}。相关研究表明，多数高校采用自有车辆、社会化租赁或两者结合的方式提供班车服务，以保障师生的基本出行需求\cite{zhang2025}。然而，传统班车模式存在班次固定、信息公开不足、缺乏个性化服务等问题，难以满足师生日益多样化的通勤需求\cite{xu2019}。

近年来，随着智慧校园建设的推进，校车及校园公交服务逐渐实现数字化管理。基于微信小程序、GPS 定位及实时数据处理的综合服务系统，可以实现班车行车状态实时跟踪、站点排队感知、乘车体验评价等功能\cite{zhu2024,yang2018}。这类系统通过整合实时数据和交互功能，不仅提升了班车调度效率，优化了班车与校园公交运行，也改善了师生乘车体验，实现通勤信息的动态共享和可视化管理，为校园交通管理提供了科学的决策依据。

在共享经济背景下，校园拼车服务成为高校出行的重要补充。调查研究显示，大学生对校园拼车具有较高接受度，主要关注安全性、路线顺路性和成本节约\cite{wei2023,zou2020,wu2017}。现有校园拼车小程序通过用户身份认证、路线匹配、实时监控等机制，实现了拼车需求的高效撮合，有效缓解了交通资源不足、出行成本高和安全隐患等问题。同时，这类应用具备可推广性和可扩展性，为高校智能交通服务提供了可行的技术基础\cite{yuan2023,sun2016,wang2020,liang2016}。

基于上述现状，本课题旨在设计与实现一个面向校园出行的综合交通平台，通过统一的数字化入口整合校园交通信息，提升出行匹配效率和用户体验，为师生提供可靠、便捷的通勤服务。这不仅回应了当前校园出行的现实需求，也为高校智慧交通管理提供可实践的技术方案。

\section{本文主要工作}

本文的主要工作是设计并实现一个面向校园场景的综合交通平台。系统以班车信息整合、拼车服务组织和用户互动支持为核心，采用微信小程序作为统一入口，结合 Spring Boot 微服务架构完成前后端协同，实现校园出行信息的集中管理与服务提供。

在功能实现方面，本文围绕用户管理、拼车服务、社区交流、聊天通知和公共交通信息五类核心业务展开设计与实现，构建了较完整的业务闭环。在算法实现方面，本文针对校园拼车场景提出了“规则筛选 + 随机森林排序”的二层匹配机制，用于提高候选订单排序的合理性与可用性。同时，本文还对系统的数据存储结构、接口组织方式和测试验证过程进行了系统整理，为平台的稳定运行和后续扩展提供支撑。

\section{论文结构安排}

本文共分为七章，各章内容安排如下：

第一章为引言，主要介绍课题的研究背景与研究意义，并对相关研究现状、本文主要工作和全文结构进行说明。

第二章为相关技术概述，主要介绍系统实现涉及的关键技术，包括系统架构、微信小程序前端技术、Spring Boot 后端开发技术以及数据库存储技术。

第三章为需求分析，对系统总体需求、功能需求、非功能需求及可行性进行分析，为后续设计提供依据。

第四章为系统设计，详细描述系统总体架构、模块设计、数据库设计与接口设计，说明各模块间的协作关系。

第五章为系统实现，介绍前端与后端的具体实现过程，并重点说明拼车匹配算法的工程实现方式。

第六章为系统测试，通过单元测试、集成测试和系统测试验证平台功能正确性、稳定性与可用性。

第七章为总结与展望，对全文工作进行总结，并结合系统实际情况提出后续改进方向。
